\section{Answers}
\label{sec:conc-answers}

\paragraph{RQ1 passes.} The best MobileNetV2 configuration reaches a melanoma F1 of
\RQoneFoneMobile{} against \RQoneFoneResnet{} for the best ResNet-50 configuration, a gap
of \RQoneGap{} against a threshold of 0.05 fixed before any run executed. The result
holds under the alternative definition of ``best'' as highest macro F1, where the gap
narrows to \RQoneAltGap{} and reverses direction. A network roughly a tenth the size
gives up nothing measurable on the class that matters most.

\paragraph{RQ2 is architecture-dependent.} Weighted cross-entropy is best for
MobileNetV2 (\RQtwoBestMobile), augmentation-only for ResNet-50 (\RQtwoBestResnet),
where both remedies made matters worse rather than better.
There is no single best imbalance strategy for this dataset; there is a best strategy per
backbone, and Chapter~8 traces the difference to optimisation stability rather than to
anything about the classes.

\paragraph{RQ3 favours MobileNetV2 on every measure.} It is the smallest at 9.07~MB
against ResNet-50's 94.30~MB, holds 10.5 times fewer parameters, and is the fastest of the
three on both CPU architectures tested.

The unanticipated part is that parameter count did not predict latency, and that the
ranking of the two heavier models was not even stable across hosts. EfficientNet-B3 holds
2.2 times fewer parameters than ResNet-50 yet runs 2.9 times slower on ARM64 and 2.0 times
faster on x86-64. Depthwise separable convolutions trade arithmetic volume for memory
traffic, and whether that trade pays depends on where the host binds.

Two things follow. Efficiency reported as parameter count, which is what the literature
usually reports, does not predict what a model costs to run. And an efficiency claim that
does not name the CPU it was measured on is not a claim at all. MobileNetV2's advantage
is the only part of the picture robust to that choice, which is what makes it the
defensible deployment recommendation.

\paragraph{The ablation does not change the recommendation.} EfficientNet-B3 produced the
highest-scoring run in the study (E11, macro F1 0.738), but at matched input resolution it
is statistically indistinguishable from MobileNetV2 in all three comparisons, and its
advantage traces to the 300-pixel input rather than to the architecture. At that
resolution it costs 8.9 times MobileNetV2's inference time on ARM64. It is the better
choice for a server-side deployment and the worse one for the offline mobile setting this
thesis is about.

\section{Contributions}
\label{sec:conc-contrib}

The first is the controlled factorial itself. Every configuration shares a split
verified by hash, an augmentation pipeline, an optimiser, a schedule and a GPU, so
differences between cells are attributable to the two factors that varied. No published
HAM10000 study identified in Chapter~2 makes that claim, and it is what allows the RQ2
finding to be stated as a property of the architectures rather than an artefact of two
papers using different setups.

The second is per-class reporting throughout. Melanoma F1 varies by a factor of two
across the six confirmatory runs, from 0.254 to 0.506, while overall accuracy varies by
0.10 and, between E1 and E2, by 0.004. Any study reporting only the aggregate would
present those two configurations as near-equivalent.

The third is the finding that imbalance remediation is not architecture-neutral. Weighted
cross-entropy helped MobileNetV2 and harmed ResNet-50 under otherwise identical
conditions. A recommendation of an imbalance strategy without reference to the backbone
is therefore incomplete.

The fourth is diagnostic rather than experimental. Chapter~8 documents that checkpoint
selection on validation loss cost up to 0.067 macro F1, concentrated on the architecture
with the noisier optimisation. That is a limitation of a protocol this thesis inherited
and followed, and identifying it is more useful to anyone repeating this work than the
headline numbers are.

The fifth is the resolution ablation. Running EfficientNet-B3 at both 224 and its native
300 pixels, with everything else fixed, separates architectural merit from input size.
Published comparisons vary the two together and therefore cannot attribute an
EfficientNet advantage to either. The answer here is that resolution contributes
materially in one of three training conditions and negligibly in the other two.

The sixth is the implementation. Every table and figure is generated from stored run
artefacts by script. A number in this document that disagreed with the experiments would
require the generating code to be wrong rather than a typo to have slipped through.

\section{What This Does Not Establish}
\label{sec:conc-limits}

Each configuration was trained once. The McNemar tests establish that two specific
trained models differ on per-image correctness; they do not establish that two
configurations differ in expectation across seeds. Every claim here should be read with
that scope.

E4 did not converge within its epoch budget. The findings are HAM10000-specific by
design. The dataset carries no skin-tone annotation, so nothing here speaks to
performance across Fitzpatrick types, which for a screening tool is a limitation of
consequence rather than of completeness.

\section{Future Work}
\label{sec:future}

\paragraph{Select checkpoints on the reported metric.} Section~\ref{sec:disc-earlystop}
shows selection on validation loss cost up to 0.067 macro F1 and did so unevenly across
architectures. Repeating the factorial with selection on validation macro F1 is a small
change that would remove a confound affecting the headline comparison. This is the first
thing to do.

\paragraph{Multiple seeds.} Three to five seeds per cell would separate genuine
differences from run-to-run variance and would let the analysis report confidence
intervals instead of point estimates. It costs a linear multiple of the compute and is
the largest single improvement available to this design.

\paragraph{A regularised re-run.} No weight decay, schedule or gradient clipping was
used. Repeating with a non-zero $L_2$ penalty and a short warmup would test whether
ResNet-50's instability is architectural, as Chapter~8 argues, or an artefact of an
under-regularised optimiser configuration. The current design cannot separate the two.

\paragraph{Matched dropout.} Adding $p = 0.2$ before ResNet-50's classifier would
quantify how much of MobileNetV2's margin comes from regularisation rather than
architecture. It is a one-line change.

\paragraph{Extend the resolution ablation to the other backbones.} The B3 pairs showed
that input resolution matters only in combination with weighted cross-entropy
(Section~\ref{sec:results-ablation}). Whether that interaction is a property of
EfficientNet or of the training condition is not answerable from one architecture.
Running MobileNetV2 and ResNet-50 at a second resolution would separate the two.

\paragraph{Skin-tone-stratified evaluation.} On a dataset carrying Fitzpatrick
annotations, or against an inferred estimate. Until this is done, no deployment
recommendation from this work should be read as applying across skin tones.

\paragraph{Qualitative error analysis.} Grad-CAM~\cite{selvaraju2017gradcam} on the
recommended configuration, and a misclassified-image review for the malignant classes,
to check the models attend to lesions rather than to acquisition artefacts.

\paragraph{On-device measurement.} Latency here is measured on a desktop CPU. A figure
from an actual mobile SoC would speak directly to the deployment claim, which desktop
CPU timing only approximates.

\section{Closing}
\label{sec:conc-closing}

The question was what a lightweight architecture costs on the classes a screening tool
cannot afford to miss. On HAM10000, under a controlled protocol, the answer is that it
costs \RQoneGap{} of melanoma F1 against a ten-fold reduction in parameters, and that
under two of three imbalance strategies the two architectures cannot be told apart on
per-image correctness at all.

The more useful result is the one that was not anticipated. The best imbalance strategy
turned out to depend on the architecture it was applied to, and the mechanism was
optimisation stability rather than class frequency. That is not visible in any study that
varies one factor at a time, and it is the kind of thing a controlled factorial exists to
find.
